\chapter{Giới thiệu}

Mạng di động đang chuyển từ phần cứng chuyên dụng sang phần mềm trên máy chủ thông thường.
Các chức năng từng chạy trên thiết bị tích hợp theo chiều dọc, chuyên biệt cố định, nay được
tái hiện thực thành chương trình trên máy tính đa mục đích. Mạng truy nhập vô tuyến (RAN),
lâu nay là phần phụ thuộc phần cứng nhất trong ngăn xếp, cũng nằm trong xu hướng đó.

Trong 5G, trạm gốc (gNB) được tách thành một đơn vị trung tâm và một hoặc nhiều đơn vị phân
tán. Đơn vị trung tâm lại được tách thành các chức năng mặt điều khiển và mặt người
dùng~\cite{3gpp38401,3gpp38300}. Open-RAN đẩy cùng ý tưởng xa hơn: RAN trở thành phần mềm
ảo hoá, đa nhà cung cấp, nối với nhau qua giao diện
mở~\cite{polese2023oran}. Mặt điều khiển thu được, gNB-CU-CP, là một chương trình lớn, có
trạng thái, tương tranh cao. Nó theo dõi mọi thiết bị đang kết nối, chạy báo hiệu
RRC~\cite{3gpp38331}, và điều phối di động cùng chuyển giao, tất cả trong giới hạn độ trễ
nghiêm ngặt.

Tải tại một site bận rộn rất cao. Một site có thể chạy hàng nghìn thủ tục quản lý kết nối và
di động mỗi giây, mỗi thủ tục có ngân sách vài
mili giây~\cite{3gpp38913}. Một mặt điều khiển triển khai có thể theo dõi hàng triệu thiết bị
hoạt động cùng lúc. Khi RAN chuyển sang triển khai đám mây, tách rời, tính đúng đắn, độ tin
cậy và thời gian của phần mềm này ngày càng quyết định chất lượng của toàn mạng. Một sự cố
hoặc một đỉnh độ trễ đuôi ở đây không phải lỗi cục bộ. Mọi thiết bị site phục vụ đều bị
ảnh hưởng~\cite{dean2013tail,barroso2017microseconds}.

Phần mềm mặt điều khiển này, hiện nay, chủ yếu được viết bằng C++. Lựa chọn đó là hợp lý:
C++ mang lại kiểm soát mức thấp, công cụ trưởng thành và hiệu năng mà hệ thống nhạy độ trễ
đòi hỏi. Nhưng hiệu năng đó đi kèm chi phí đã được biết. C++ đặt gánh nặng quản lý bộ nhớ và
tương tranh lên lập trình viên, và hàng thập kỷ kinh nghiệm cho thấy gánh nặng này thường
xuyên bị xử lý sai, sinh ra lỗi use-after-free, double-free, tranh chấp dữ liệu và các lỗi
liên quan, tốn kém để tìm, nguy hiểm khi vận hành, và chiếm phần lớn lỗ hổng bảo mật trong
thực tiễn công nghiệp~\cite{msrc2019memory,google2022android}. Rust đã nổi lên như ngôn ngữ
hệ thống đại chúng đầu tiên hứa hẹn thoát khỏi đánh đổi này: an toàn bộ nhớ và không tranh
chấp dữ liệu được thực thi lúc biên dịch, với hiệu năng dự kiến cạnh tranh với
C++~\cite{klabnik2023rust,jung2018rustbelt}. Liệu lời hứa đó có đúng với mặt điều khiển cấp
viễn thông, nhạy độ trễ, nơi khả năng dự đoán khi quá tải quan trọng ngang hiệu năng thuần,
vẫn chưa được hiểu rõ. Luận văn này, thực hiện tại Nokia, trả lời câu hỏi trực tiếp: phát
triển hai bản hiện thực tương đương về chức năng của một phân hệ RAN đại diện, một dịch vụ
Quản lý Ngữ cảnh UE mô hình hoá theo một thành phần gNB-CU-CP thật của Nokia, bằng C++ hiện
đại và Rust, rồi so sánh chúng dưới khối lượng báo hiệu tổng hợp sát thực tế.

Ba thuật ngữ lặp lại xuyên suốt. \emph{gNB-CU-CP} là phần mềm mặt điều khiển mà công trình
này xoay quanh. \emph{Quản lý Ngữ cảnh UE (UECM)} là phân hệ cụ thể được tái hiện thực và đo
lường. \emph{Độ trễ đuôi}, các phân vị p99 và p99.9 thay vì trung bình, là thuộc tính quyết
định liệu hệ thống như vậy có đáp ứng hạn chót hay
không~\cite{dean2013tail,barroso2017microseconds}. Phụ lục~\ref{appendix:terms} cố định các
thuật ngữ này cùng phần còn lại của từ vựng miền cho độc giả ngoài ngành viễn thông.

\section{Động lực và bối cảnh công nghiệp}\label{sec:motivation}
Ba xu hướng khiến phép so sánh này đúng thời điểm. Xu hướng thứ nhất là tiến hoá đám mây của
chính RAN. Đưa các chức năng mặt điều khiển lên máy chủ đa mục đích làm tăng, chứ không giảm,
đòi hỏi về chất lượng phần mềm: phần mềm sống lâu hơn, liên tục tiến hoá, chạy với mức tương
tranh cao hơn, và phải đáp ứng hạn chót mà không còn tính xác định mà phần cứng chuyên dụng
từng cung cấp. Xu hướng này được đẩy nhanh bởi tách rời Open-RAN, biến trạm gốc thành phần
mềm ảo hoá, đa nhà cung cấp, nối qua giao diện
mở~\cite{polese2023oran}, nên mặt điều khiển mà nhà cung cấp giao hàng ngày càng cạnh tranh
bằng chất lượng kỹ thuật phần mềm (độ tin cậy, thời gian dự đoán được, bảo mật, và mức độ dễ
tiến hoá an toàn một mã nguồn lớn) thay vì phần cứng chuyên dụng. Trong bối cảnh đó, một hồi
quy độ trễ hoặc một lỗi hỏng bộ nhớ không còn là lỗi cục bộ mà là rủi ro tin cậy và bảo mật
toàn mạng, và câu hỏi chọn ngôn ngữ không còn học thuật mà trở thành quyết định kỹ thuật cụ
thể.

Xu hướng thứ hai là AI chuyển từ công cụ tinh chỉnh mạng sang khối lượng công việc mà mạng
được xây để vận chuyển, và thế hệ kế tiếp đang được đặc tả theo hướng đó ngay từ đầu.
IMT-2030, khung ITU cho 6G, liệt kê trí tuệ nhân tạo và truyền thông trong sáu kịch bản sử
dụng và trí tuệ phổ biến trong bốn nguyên tắc thiết
kế~\cite{itu2023imt2030}. Kiến trúc hình thành quanh đó là AI-RAN, trong đó viễn thông và
khối lượng AI được điều phối chung trên hạ tầng gia tốc dùng chung, nên các site chạy vô
tuyến cũng lưu trữ nền tảng AI biên phân tán~\cite{polese2026airan}. Liên minh AI-RAN nêu
thẳng yêu cầu vận hành: một hạ tầng chung phải chạy đồng thời cả hai loại khối lượng, và
thỏa thuận mức dịch vụ của chúng có thể khác
nhau~\cite{airan2025orchestrator}. Điều đó thay đổi những gì lựa chọn ngôn ngữ phải mang lại
cho mặt điều khiển. Phần mềm của nó không còn chiếm trọn máy chủ. Nó tranh CPU, bộ nhớ và
cache với công việc suy luận có mẫu đến nó không kiểm soát, và đồng ở việc chiếm dụng làm
giảm tải mà mặt điều khiển bão hoà. Bão hoà là nơi hai ngôn ngữ tách biệt rõ nhất trong các
đo lường của Chương~\ref{chapter:eval}. Dấu chân bộ nhớ và hành vi sau điểm đó đang trở thành
tiêu chí lựa chọn riêng, và không tiêu chí nào được quyết định bởi kiến trúc phía trên mã
nguồn.

Xu hướng thứ ba là chuyển dịch toàn ngành sang ngôn ngữ an toàn bộ nhớ. Các tổ chức phần mềm
lớn báo cáo phần lớn lỗ hổng bảo mật nghiêm trọng bắt nguồn từ lỗi an toàn bộ nhớ, và việc
áp dụng ngôn ngữ an toàn bộ nhớ giảm chúng đo
được~\cite{msrc2019memory,google2022android,vanderstoep2025android,google2024securebydesign}.
Luận điểm không còn chỉ dựa vào báo cáo nhà cung cấp: một nhóm gồm học thuật và công nghiệp
đã lập luận rằng an toàn bộ nhớ giờ cần thuật ngữ và yêu cầu chuẩn hoá nếu muốn đặc tả và
mua sắm~\cite{watson2025memorysafety}. Mối quan tâm đã lên tầm chính sách quốc gia: cơ quan
an ninh mạng hiện khuyến khích nhà sản xuất phần mềm công bố ``lộ trình an toàn bộ nhớ'' và
chuyển khỏi ngôn ngữ không an toàn bộ nhớ như C và
C++~\cite{cisa2023memorysafe,oncd2024building}. Trong thị trường mà phân hệ này được bán
vào, vấn đề cũng trở thành quy định. Đạo luật Khả năng phục hồi mạng EU đặt nghĩa vụ thiết kế
an toàn và xử lý lỗ hổng trong vòng đời lên nhà sản xuất sản phẩm có yếu tố số, với nghĩa vụ
báo cáo lỗ hổng từ tháng 9 năm 2026 và phần lớn phần còn lại từ tháng 12 năm
2027~\cite{eu2024cra}. Văn bản không nêu ngôn ngữ lập trình, và không nên đọc nó như bắt buộc
một ngôn ngữ. Điều nó thay đổi là các lớp lỗi lặp lại của ngôn ngữ không an toàn bộ nhớ trở
thành rủi ro tuân thủ cũng như rủi ro kỹ thuật. Rust đặc biệt đã chuyển từ ngôn ngữ ngách
vào nhân hệ điều hành, hạ tầng đám mây, và ngày càng vào viễn thông. Siêu phần ảo hoá sản
xuất như AWS Firecracker, nền tảng của nền tảng serverless, đã được viết bằng
Rust~\cite{agache2020firecracker}. Song cơ sở bằng chứng về mức phù hợp của Rust chủ yếu dựa
trên phần mềm đa mục đích và benchmark hiệu năng trường hợp trung bình. Câu hỏi cụ thể mà
nhà cung cấp viễn thông phải trả lời, liệu Rust có đáp ứng yêu cầu độ trễ đuôi và khả năng
dự đoán của mặt điều khiển RAN, liệu nó có thực sự ngăn các lớp lỗi ám ảnh mã C++ hiện có,
và chi phí áp dụng tăng dần trong mã nguồn C++ lớn là bao nhiêu, chưa được trả lời bằng
bằng chứng trực tiếp, có kiểm soát. Cung cấp bằng chứng như vậy cho một phân hệ đại diện là
động lực của công trình này.

\paragraph{Bối cảnh công nghiệp.}
Luận văn này được thực hiện tại Nokia, một trong số ít nhà cung cấp xây dựng, duy trì và tiến
hoá phần mềm RAN sản xuất từ đầu đến cuối. Thị trường đó tập trung rất ít tay: năm nhà cung
cấp lớn nhất chiếm 96\% doanh thu RAN toàn cầu trong ba quý đầu năm 2025, với chỉ số
Herfindahl--Hirschman ở mức cao nhất trong mười
năm~\cite{delloro2025ran3q,delloro2026ran}. Phân hệ nghiên cứu, Quản lý Ngữ cảnh UE (UECM),
được mô hình hoá theo một thành phần mặt điều khiển gNB-CU-CP thật chứ không phải ví dụ sách
giáo khoa.

gNB-CU-CP không phải thành phần phụ. Trong kiến trúc tách rời 3GPP, nó là nửa mặt điều khiển
của đơn vị trung tâm, lưu trữ RRC và quản lý vòng đời ngữ cảnh
UE~\cite{3gpp38401,3gpp38331}, dưới ngân sách thủ tục cấp mili giây mà hệ thống 5G đặt lên
báo hiệu điều khiển~\cite{3gpp38913} và ở tốc độ sự kiện cùng số thiết bị kết nối mà mạng
triển khai duy trì. Mục~\ref{sec:ranarch} đặt nó trong kiến trúc rộng hơn. Vì được viết bằng
C++, nó thừa hưởng cả hiệu năng lẫn rủi ro an toàn bộ nhớ và tương tranh của ngôn ngữ đó. Đó
chính là nơi ba xu hướng trên hội tụ, và phần mềm loại này được giao hàng, tiến hoá và vận
hành trên mạng nhà mạng thật, nên quyết định chọn ngôn ngữ là mối quan tâm sản xuất chứ
không phải giả định.

UECM là đối tượng nghiên cứu có chủ đích trong bối cảnh đó, vì các lý do Mục~\ref{sec:uecm}
nêu: đủ nhỏ để tái hiện thực trung thực và lý giải toàn bộ, song vẫn mang các thuộc tính
khiến mã mặt điều khiển khó. Hiện thực hoá hoàn toàn theo một đặc tả công khai, đã đóng băng,
chứ không theo mã độc quyền, để phép so sánh có thể tái lập bởi độc giả ngoài Nokia.

\paragraph{Một quyết định mang tầm một thế hệ.}
Lựa chọn ngôn ngữ trong phần mềm loại này không thường xuyên được xem xét lại, và nó sống
lâu hơn thế hệ mà nó được đưa ra. Thế hệ kế tiếp đang được đặc tả ngay bây giờ: 3GPP đã
phân bổ Release 20 cho nghiên cứu 6G và Release 21 cho công việc chuẩn hoá, theo lịch ITU
yêu cầu đề xuất công nghệ đầu năm 2029 và nộp đặc tả giữa
2030~\cite{3gpp2024road6g}. Công việc đó bao gồm rõ ràng di chuyển từ và tương tác với 5G, nên
mặt điều khiển viết cho thế hệ hiện tại được duy trì song song thế hệ kế tiếp thay vì nghỉ
hưu khi thế hệ đó tới. Mục tiêu thế hệ sau cũng không nới lỏng động lực của phép so sánh
này. IMT-2030 đặt bảo mật và khả năng phục hồi là một trong bốn nguyên tắc thiết kế, kỳ vọng
hệ thống an toàn ngay từ thiết kế, định nghĩa khả năng phục hồi là khả năng tiếp tục vận hành
đúng trong và sau sự kiện gây rối, và mang kịch bản độ trễ thấp siêu tin cậy về ở dạng củng
cố~\cite{itu2023imt2030}. Thoát khỏi cả lớp lỗi và hành vi dự đoán được khi bị nhiễu đang
được ghi vào mục tiêu thế hệ kế tiếp thay vì giao cho người hiện thực.

Các mốc chính sách nêu trên nằm trong cùng cửa sổ. CISA và FBI xếp phát triển dòng sản phẩm
mới bằng ngôn ngữ không an toàn bộ nhớ là thực hành xấu khi sản phẩm phục vụ hạ tầng quan
trọng, phạm vi mà hạ tầng mạng di động nằm trọn trong đó, và họ miễn sản phẩm hiện có khỏi
kỳ vọng lộ trình an toàn bộ nhớ chỉ khi thời hạn hỗ trợ công bố kết thúc trước tháng 1 năm
2030~\cite{cisa2025badpractices}. Nhà cung cấp hoãn câu hỏi ngôn ngữ đến khi thế hệ kế tiếp
được đặc tả sẽ không tránh được việc quyết định. Họ sẽ quyết định mặc định, bằng ngôn ngữ
hiện hành, cho thêm một thập kỷ triển khai.

Điều luận văn này đóng góp cho quyết định đó có giới hạn, và giới hạn đáng nêu rõ. Các đo
lường ở đây lấy trên phân hệ mô hình hoá theo thành phần thế hệ hiện tại và đánh giá theo
yêu cầu đã công bố. Không có tuyên bố về yêu cầu chưa được viết. Điều vượt ranh giới thế hệ
không phải một con số độ trễ mà chính phép so sánh cấp ngôn ngữ: lớp lỗi nào trình biên dịch
có thể loại trừ, mô hình an toàn tốn bao nhiêu, và áp dụng tăng dần trong mã nguồn C++ lớn
gồm những gì. Các câu hỏi đó có cùng câu trả lời bất kể đặc tả phía trên gọi thế hệ nào, và
rẻ nhất khi trả lời trước khi bắt đầu mã nguồn kế tiếp chứ không phải sau.

\paragraph{Một kiểu lỗi cụ thể.}
Để làm rõ mức độ, hãy xét một cơn bão di động khi nhiều UE chuyển giao cùng lúc. Một con trỏ
xử lý sai, ngữ cảnh UE được giải phóng khi thả trong khi thông điệp báo hiệu đang bay vẫn
tham chiếu tới nó, sinh ra use-after-free. Trong C++ mã này biên dịch sạch và có thể vượt
kiểm thử thường, chỉ để hỏng trạng thái UE không liên quan hoặc làm sập luồng worker đúng
điều kiện tương tranh cao mà hệ thống tồn tại để xử lý. Cùng một tranh chấp tiềm ẩn cũng có
thể xuất hiện thành đỉnh độ trễ đuôi khó đoán thay vì sập hẳn. Trong Rust an toàn cùng thiết
kế không biên dịch: sở hữu và mượn từ chối tham chiếu treo lúc build. Kịch bản này nắm lý do
luận văn đo cả \emph{an toàn} (lỗi nào được ngăn, và khi nào) lẫn \emph{khả năng dự đoán}
(độ trễ đuôi khi quá tải) thay vì chỉ thông lượng trung bình.

\section{Phát biểu vấn đề}
Hai thuộc tính của miền này khiến lựa chọn ngôn ngữ có hệ quả. Thứ nhất, quản lý bộ nhớ và
tương tranh thủ công trong C++ để phần mềm mắc các lỗi như use-after-free, double-free, tranh
chấp dữ liệu và vô hiệu hoá iterator, im lặng lúc biên dịch và tốn kém để tìm sửa ở quy mô
viễn thông. Thứ hai, điều quyết định liệu hạn chót RAN được đáp ứng không phải thông lượng
trung bình mà hành vi trường hợp xấu nhất và đuôi dưới tải và quá
tải~\cite{dean2013tail}, thuộc tính dễ hồi quy và khó đo đúng. Rust hứa loại bỏ lớp lỗi thứ
nhất lúc biên dịch trong khi tuyên bố hiệu năng cạnh tranh ở lớp thứ hai, song tuyên bố đó
chủ yếu dựa trên benchmark đa mục đích thay vì khối lượng viễn thông có trạng thái, tương
tranh, nhạy độ trễ.

Gộp các mạch lại, vấn đề cụ thể mà luận văn này giải quyết như sau. Nhà cung cấp viễn thông
chịu áp lực ngày càng lớn, cả kỹ thuật lẫn quy định, để áp dụng ngôn ngữ an toàn bộ nhớ, song
với mặt điều khiển RAN họ thiếu bằng chứng \emph{trực tiếp, có kiểm soát} để căn cứ quyết
định đó. Chưa được thiết lập liệu bản hiện thực Rust của phân hệ mặt điều khiển có trạng thái,
tương tranh có khớp C++ hiện đại về \emph{độ trễ đuôi và khả năng dự đoán khi quá tải} (trái
với thông lượng trung bình), liệu đảm bảo lúc biên dịch của Rust có thực sự loại bỏ các lớp
lỗi cụ thể ảnh hưởng mã C++ hiện có, và chi phí áp dụng tăng dần \emph{là bao nhiêu} trong
mã nguồn C++ lớn. Bằng chứng sẵn có chủ yếu đa mục đích và trường hợp
trung bình~\cite{balasubramanian2017rust,qin2020rust}, để ba câu hỏi (khả năng dự đoán, phạm
vi lỗi bị ngăn, và chi phí di chuyển) mở với một phân hệ đại diện. Trả lời chúng trong điều
kiện có kiểm soát, tái lập được là vấn đề mà luận văn này đặt ra để giải quyết.

\section{Câu hỏi nghiên cứu và giả thuyết}\label{sec:rqs}
Luận văn được tổ chức quanh bốn câu hỏi nghiên cứu, sắp theo thứ tự ưu tiên khả năng dự đoán
và di chuyển trước hiệu năng trung bình thuần:

\begin{description}
  \item[RQ1] Rust và C++ hiện đại so sánh thế nào về khả năng dự đoán độ trễ (trung vị, độ
        trễ đuôi p99 và p99.9, và giá trị cực đại) cho khối lượng mặt điều khiển RAN dưới tải
        tăng dần và quá tải?
  \item[RQ2] Lớp lỗi an toàn bộ nhớ và tương tranh nào xảy ra ở C++ bị loại bỏ lúc biên dịch
        trong bản hiện thực Rust, và các lỗi còn lại được phát hiện ở giai đoạn nào với chi phí
        bao nhiêu?
  \item[RQ3] Chi phí, nếu có, của mô hình an toàn Rust (kiểm tra biên, sao chép do sở hữu, chi
        phí trừu tượng hoá) lên thông lượng, sử dụng CPU và hành vi bộ nhớ là gì?
  \item[RQ4] Hệ quả thực tiễn (công sức, khả năng tương tác FFI và rủi ro) của việc áp dụng
        Rust tăng dần trong hệ sinh thái phần mềm RAN C++ sẵn có là gì?
\end{description}

Các câu hỏi được tiếp cận với ba kỳ vọng trước, nêu ở đây để phần đánh giá có thể xác nhận
hoặc bác bỏ thay vì chỉ mô tả kết quả sau sự kiện. \textbf{H1:} với khối lượng này hai ngôn
ngữ tương đương hiệu năng trong vùng vận hành liên quan dưới bão hoà, mọi khác biệt chỉ nằm
ở đuôi xa. \textbf{H2:} Rust an toàn ngăn lúc biên dịch các lớp lỗi an toàn bộ nhớ và tương
tranh mà C++ giao cho phát hiện nỗ lực tốt nhất, sau sự kiện. \textbf{H3:} chi phí chạy của
tương tác C++--Rust là không đáng kể so với chi phí các thao tác được gọi, nên chi phí thật
của áp dụng tăng dần là công sức kỹ thuật chứ không phải hiệu năng. H2 được phát biểu để có
thể sai. Danh mục lỗi phía sau RQ2 trải cả lớp Rust an toàn dự kiến ngăn và lớp dự kiến không
ngăn, để nghiên cứu định vị ranh giới đảm bảo thay vì chỉ xác nhận. Các giả thuyết được xem
lại theo bằng chứng ở Mục~\ref{sec:hypotheses}.

\section{Đóng góp, kết quả và hạn chế}
Luận văn có sáu đóng góp. Hai là hiện vật: một khung benchmark tái lập được cho khối lượng
mặt điều khiển viễn thông, và một đặc tả chung đã đóng băng với hai bản hiện thực đã kiểm thử
tương đương của phân hệ mặt điều khiển RAN, một bằng C++ và một bằng Rust. Ba là kết quả.
Thứ nhất là so sánh độ trễ và tài nguyên chia tại điểm gối bão hoà, nơi thứ tự giữa hai ngôn
ngữ đảo ngược. Thứ hai là phân loại lớp lỗi từ tiêm lỗi có kiểm soát đánh dấu nơi đảm bảo lúc
biên dịch dừng lại, theo cả hai hướng. Thứ ba là giới hạn trên những gì phép so sánh dạng này
có thể quy cho ngôn ngữ, thu được bằng cách phân rã một thao tác đầu-cuối thành công việc
ngôn ngữ thực hiện và bàn giao chéo luồng quanh nó. Cuối cùng là tập hướng dẫn kỹ thuật và di
chuyển cho chọn ngôn ngữ trong tiến hoá RAN đám mây (Mục~\ref{sec:guidelines}). Cả sự đảo
ngược lẫn giới hạn quy kết đều không được dự đoán trước. Cả hai được tìm thấy khi đọc kết quả
mà nghiên cứu dự định báo cáo theo cách khác, và Mục~\ref{sec:hypotheses} ghi chi phí của cái
đầu đối với H1.

Các phát hiện chia ba nhánh. Về \emph{an toàn} (RQ2), mọi lớp lỗi an toàn bộ nhớ và tranh chấp
dữ liệu được tiêm đều bị ngăn lúc biên dịch trong Rust an toàn. Cùng các lỗi đó đều biên dịch
trong C++, và mỗi lỗi chỉ bị bắt, nếu có, lúc chạy và chỉ dưới bộ dò lỗi đúng, hoặc với một
lớp thì không có bộ dò nào. Bảy lớp ranh giới được chọn có chủ đích, và không ngôn ngữ nào
ngăn chúng. Với hai lớp trong số đó công cụ C++ đi trước, và với một lớp lỗi mà Rust an toàn
loại trừ quay lại ngay khi chương trình viết \texttt{unsafe}. Về \emph{di chuyển} (RQ4), đường
ghép FFI nhỏ và giữ lõi không cần \texttt{unsafe}. Đo so với cùng thao tác gọi bản địa, chi
phí tối đa $\FfiOpBoundPct\%$ của một thao tác, nên chi phí áp dụng thật là thiết kế API chứ
không phải chi phí chạy. Ba chi phí ngoài đường ghép được đo vì cùng lý do, và hai trong số
đó chống lại quan niệm thông thường: phần lớn chi phí build thêm của Rust là biên dịch mã bên
thứ ba chứ không phải công việc của chính trình biên dịch, và thứ tự kích thước nhị phân đảo
ngược khi hai chuỗi công cụ được giữ cùng một chính sách liên kết. Cả hai kết quả này là cuối
cùng.

Về \emph{hiệu năng} (RQ1, RQ3) câu trả lời chia hai lần, và vẽ đúng cả hai phân chia là nơi
phần lớn phân tích đi vào. Phân chia thứ nhất tách những gì phép quét đo khỏi những gì nó có
thể quy kết. Một thao tác đầu-cuối trong thiết kế này là $\BudgetCoreShareMax\%$ công việc
ngôn ngữ và bàn giao chéo luồng chiếm phần còn lại, nên so sánh cặp khớp trên đường đó giới
hạn thiết kế chứ không phải ngôn ngữ. Giới hạn đó đúng bất kể phép so sánh như vậy ra sao, và
chạm mọi nghiên cứu báo cáo trung vị đầu-cuối mà không đo tỷ lệ phần. Đo tách khỏi bàn giao,
thao tác nhanh hơn $\BudgetCoreRatio$ lần trong Rust, điều mà phép quét không thấy.

Phân chia thứ hai là theo chế độ. Dưới bão hoà hai bản hiện thực tương đương đầu-cuối về hành
vi trung tâm, với biên độ không triển khai nào nhận ra. Kiểm định cặp vẫn phân giải khác biệt
trung vị $\PFiftyDiffAbsMax\,\mu$s nghiêng về Rust dưới cả hai bộ cấp phát, với khoảng tin cậy
chạm $\PFiftyDiffCIMax\,\mu$s, nên tương đương là thực tiễn chứ không phải thống kê. Đuôi xa
cùng chế độ không thể kết luận trên máy chủ sẵn có. Qua bão hoà thứ tự đảo ngược. Bản C++ duy
trì $\OverloadThrRatio$ lần thông lượng của bản Rust trong \OverloadThrAgree{} lần lặp đan xen
dưới cả hai bộ cấp phát, và dùng nhiều CPU hơn cho mỗi thao tác hoàn thành. Về tài nguyên tiến
trình Rust giữ thêm vài MiB bộ nhớ cư trú, lại dưới cả hai bộ cấp phát, và nhỏ hơn một bậc
độ lớn so với chi phí của chính lựa chọn bộ cấp phát.

Bốn hạn chế bao quanh tất cả điều này: khối lượng tổng hợp, một phân hệ đơn giản hoá, thiên
lệch hiện thực một tác giả, và máy chủ đo ảo hoá. Chương~\ref{chapter:discussion} xử lý từng
hạn chế chi tiết.

Phạm vi được vẽ có chủ đích để bảo vệ tính hợp lệ nội bộ của phép so sánh ngôn ngữ. Trong
phạm vi là chức năng Quản lý Ngữ cảnh UE của Chương~\ref{chapter:context}, thiết kế tương
tranh không chia sẻ, phép quét tải đến quá tải, danh mục tiêm lỗi, và nguyên mẫu FFI. Ngoài
phạm vi rõ ràng là toàn hệ gNB và các miền mặt điều khiển còn lại, xử lý PHY/L2 và mặt người
dùng nặng dữ liệu, trace lưu lượng thật (thay vì khối lượng tổng hợp có căn cứ đặc tả), các
lựa chọn ngôn ngữ có quản lý (loại trừ vì lý do độ trễ đuôi ở Chương~\ref{chapter:background}),
và triển khai sản xuất cùng mối quan tâm vận hành như nâng cấp, quan sát được và đa thuê bao.
Các loại trừ này đổi tính hợp lệ ngoại bộ lấy điều kiện có kiểm soát mà phép so sánh ngôn
ngữ công bằng cần, và mỗi loại trừ được xem lại như hạn chế ở
Chương~\ref{chapter:discussion}.

\section{Ý nghĩa và đối tượng hưởng lợi}
Kết quả hướng tới hành động trực tiếp. Với nhà cung cấp RAN và nhà vận hành mạng đám mây,
luận văn cung cấp bằng chứng có kiểm soát cho câu hỏi chặn mọi quyết định ngôn ngữ. Câu hỏi
đó không phải ``Rust có nhanh trung bình không?'' mà ``Rust có giữ độ trễ đuôi và khả năng
dự đoán khi quá tải, và áp dụng tăng dần tốn bao nhiêu?'', được trả lời cho phân hệ đại diện
mã mặt điều khiển sản xuất. Với thảo luận tiêu chuẩn và chính sách quanh lộ trình an toàn bộ
nhớ~\cite{cisa2023memorysafe,oncd2024building}, công trình neo một mandat chung vào một điểm dữ
liệu viễn thông cụ thể. Với cộng đồng nghiên cứu, đặc tả đóng băng, các bản hiện thực đã kiểm
thử tương đương, khung benchmark tái lập được và phân loại lớp lỗi tạo nền tái sử dụng cho
nghiên cứu ngôn ngữ và hệ thống tiếp theo, và phân rã độ trễ cho nghiên cứu như vậy cách kiểm
tra những gì đo lường của chính nó có thể quy kết. Với kỹ sư thực hành, danh mục lỗi và nguyên
mẫu di chuyển đồng thời là ví dụ làm việc về những gì đảm bảo của Rust mua và không mua trong
miền này.
